% Drawing 129/06-1 -- folder 129, batch 06, archive page 106 (PDF page 107).
% Transcription: dessins redécrits, non reproduits
% Kind: filtrations drawn as graduated segments. Method 2 (schematic redraw from the facsimile). Status: ok.
% Notes: Two graduated segments with their labels written vertically along the ticks: below, A > fA > f^2A > ... > f^{n-j}A > f^{n-j-1}A > ... > f^{n-1}A > f^nA = 0, with a brace under the first steps (A to f^{n-j}A) ending in a short stroke; above, ... f^jA > f^{j-1}A > ... > f^{n-1}A, with a brace under it. The upper labels follow the transcription (f^jA, f^{j-1}A, f^{n-1}A); a fourth label at the right end of the upper segment is overwritten and left out. A small graduated sketch further left on the page (near 'f^2A', around x 0.26-0.31, y 0.71-0.77) is not covered by the mark and is not drawn.
% Coordinates: 1 unit = 1 mm = 1/100 of the crop width (bbox in page fractions: [0.38, 0.62, 0.78, 0.82]).
% Redrawn by Claude (claude-opus-5-5), issue #73, 2026-11-07; not checked by a human; not compiled (no TeX in the session).
% Rights: reproduces part of the work -- behind RENDER_TRANSCRIPTS, never published by itself.
\documentclass[tikz,border=3pt]{standalone}
\usepackage{amsmath,amssymb}
\usetikzlibrary{patterns,arrows.meta}
\tikzset{pencil/.style={black!45}, struck/.style={black!35}, mark/.style={circle,draw,inner sep=1.3pt}, pt/.style={circle,fill,inner sep=0.9pt}, opt/.style={circle,draw,fill=white,inner sep=1.1pt}}
\begin{document}
\begin{tikzpicture}[x=1mm,y=1mm,line cap=round,line join=round,font=\scriptsize]
\node[inner sep=1pt] at (23,-18.3) {\tiny $\cdots$};
\draw (31,-18.5) -- (93,-17);
\draw[very thick] (31,-16.3) -- (31,-20.7);
\draw (42.5,-16.7) -- (42.5,-19.7);
\draw (86.5,-15.7) -- (86.5,-18.7);
\draw (93,-14.8) -- (93,-19.2);
\node[inner sep=1pt,rotate=90] at (30,-11) {\tiny $f^jA$};
\node[inner sep=1pt,rotate=90] at (41.5,-7) {\tiny $f^{j-1}A$};
\node[inner sep=1pt,rotate=90] at (84.5,-8) {\tiny $f^{n-1}A$};
\draw plot[smooth, tension=0.55] coordinates {(34,-21.5) (45,-23) (58,-23.2) (61,-25.5) (64,-23) (78,-22.5) (89,-21.5)};
\draw (8.5,-58.5) -- (73.5,-53.3);
\draw (8.5,-55.9) -- (8.5,-61.1);
\draw (20,-55.6) -- (20,-60);
\draw (31.5,-55.8) -- (31.5,-58.6);
\draw (41,-55.1) -- (41,-57.9);
\draw (52,-53.5) -- (52,-57.9);
\draw (62.5,-52.8) -- (62.5,-57.2);
\draw (63.5,-52.8) -- (63.5,-57.2);
\draw (73.5,-51.1) -- (73.5,-55.5);
\node[inner sep=1pt] at (78.5,-52.5) {\tiny $\cdots$};
\draw (82.5,-52.6) -- (91,-52);
\draw (82.5,-50.4) -- (82.5,-54.8);
\draw (91,-49.2) -- (91,-54.8);
\node[inner sep=1pt,rotate=90] at (9.5,-52) {\tiny $A$};
\node[inner sep=1pt,rotate=90] at (20.5,-50.5) {\tiny $fA$};
\node[inner sep=1pt,rotate=90] at (31.5,-48.5) {\tiny $f^2A$};
\node[inner sep=1pt,rotate=90] at (51.5,-44.5) {\tiny $f^{n-j}A$};
\node[inner sep=1pt,rotate=90] at (62,-42) {\tiny $f^{n-j-1}A$};
\node[inner sep=1pt,rotate=90] at (81.5,-41.5) {\tiny $f^{n-1}A$};
\node[inner sep=1pt,rotate=90] at (91.5,-38.5) {\tiny $f^nA=0$};
\draw plot[smooth, tension=0.55] coordinates {(12,-63) (20,-63.5) (27,-63.5) (29.5,-66) (32,-62.5) (42,-61.5) (49.5,-60.5)};
\draw (29.5,-66) -- (29.5,-70.5);
\end{tikzpicture}
\end{document}
